\SourceLeaf{243}%
until I have the leisure to take up the subject again and am more at ease reading Russian books than I am now. If I am not mistaken, a paper by p. \textsc{Lobachevsky}, perhaps a translation from the \gausscyr{Записки}, is also in \textsc{Crelle}'s journal, but I cannot check that at the moment. \mbox{.\ .\ .\ .\ .\ .\ .}

\par\bigskip\Needspace{6\baselineskip}

\begin{center}\textsc{Gauss} to \textsc{Gerling}. Göttingen, 8 February 1844.\end{center}

My dearest friend, you may not be displeased if I add one or two items to the bibliographical notes I gave you in my last letter. I cannot at present go further into the subject itself.

\textsc{Lobachevsky}'s paper in \textsc{Crelle}'s journal is in vol. 17, p. 295 ff. I find that it is only a free translation of the Russian paper [\gausscyr{Воображаемая геометрія}] in the 1835 Learned Writings of the University of Kazan, where one will run up against the same difficulties as in the German paper. In the latter, on p. 296, line 10, you run up against the words \textls{J’ai démontré} \emph{etc.}, which are of little use to a reader who has only this paper. The same remark applies at the top of p. 303 to \textls{J’ai prouvé ailleurs} \emph{etc.} The earlier paper to which this \textls{seems} to refer is probably the one cited in a note to the Russian paper mentioned above as appearing under the title Über die Anfänge oder Principe der Geometrie in the \gausscyr{Казанскомъ Вѣстникѣ} (Kazan Messenger) for 1829 and 1830. It is also remarked there that a very wounding criticism of this paper appeared in No. 41 of another Russian journal, »Sohn des Vaterlandes« \gausscyr{Сынъ Отечества}, of 1834, published, I suppose, in Petersburg. \textsc{Lobachevsky} sent in a reply to the criticism, but by the beginning of 1835 it had yet to be accepted.

These bibliographical notes are of little help, of course, since a copy of the Kazan Messenger for\space%
